\documentclass{article}

\usepackage{float}
\restylefloat{table}

\usepackage{booktabs}
\usepackage{tabularx}

\title{Team Productivity: POC\\\progname}

\author{\authname}

\date{}

\input{../Comments.text}
\input{../Common.text}

\begin{document}

\maketitle

This document summarizes the contributions of each team member up to the POC
Demo.  The time period of interest is the time between the beginning of the term
and the POC demo.

\wss{Please do not delete sections or modify the table formats.  Standardization
helps the instructors and TAs to review these documents. You may add sections to
the report.}

\wss{Remove the instructor guidelines/comments in the final version.}

\wss{Fill in the requested information, like project url, TA name, supervisor name, etc.}

\section{Demo Plans}

% \wss{What is your feasibility risk?}

% \wss{What will you be demonstrating}

\paragraph{Feasibility risks.} As set out in Section 9 of the Development
Plan, FitFolio depends on two outside systems that no team member had used
before. If either cannot be made to work, the app's core loop of earning
currency through activity and investing it cannot be built.

\begin{enumerate}
\item \textbf{Reading step data from the phone's health platform.} The app
  reads steps through Google Health Connect on Android and Apple HealthKit on
  iOS, using two third-party React Native libraries:
  \begin{itemize}
  \item \texttt{react-native-health-connect}
  \item \texttt{@kingstinct/react-native-healthkit}
  \end{itemize}
  If steps cannot be read reliably, the app cannot award currency fairly.
\item \textbf{Valuing portfolios with real market prices.} Portfolio values
  must follow real prices from Finnhub, the market data provider, while
  staying within the request limits of its free plan.
\end{enumerate}

\paragraph{What we will demonstrate.} The demo shows one path through the core
loop, on a physical phone, as planned in the Development Plan:

\begin{enumerate}
\item \textbf{Read today's steps} (Risk 1). After the user grants permission,
  the app reads the day's step count and displays it. This succeeds if the
  count matches the phone's own health app and goes up after a team member
  walks and reopens the app. Both platforms are implemented
  (\texttt{src/frontend/app/services/health/steps.ts}), so the iPhone, a
  stretch goal in the Development Plan, is now part of the demo.
\item \textbf{Turn steps into virtual currency} (connects Risks 1 and 2). The
  app sends the day's total to the server
  (\texttt{PUT /v1/me/activity/daily-totals}), which awards one credit per 100
  steps, up to 300 credits a day, and never takes currency back if a later
  total is lower. The new balance is shown (\texttt{GET /v1/me/balance}).
\item \textbf{Buy a real asset with that currency} (Risk 2). The user buys
  whole shares at the current market price (\texttt{POST /v1/me/trades}). The
  server refuses an order the balance cannot cover, and refuses all orders
  while the market is closed, so the demo will be held during market hours.
\item \textbf{Show the portfolio following the market} (Risk 2). While the
  market is open, the server requests each supported asset's price from
  Finnhub once a minute and shares it across all users.
  \texttt{GET /v1/me/portfolio} values the holdings at the latest price, and
  \texttt{GET /v1/market/usage} shows how many calls were made in the last
  minute and today, to show that this interval fits within the free plan.
\end{enumerate}

\section{Team Meeting Attendance}

\begin{table}[H]
\centering
\begin{tabular}{ll}
\toprule
\textbf{Student} & \textbf{Meetings}\\
\midrule
Total & 3\\
Davis Lenover & 3\\
Preston Grewal & 3\\
Angad Chhabra & 2\\
Farid Bastoros & 3\\
Zaid Alsaadi & 2\\
\bottomrule
\end{tabular}
\end{table}

Sections 2 to 6 cover September 8, 2026, when classes started, to November 19,
2026. The counts here come from the attendance list in each team meeting
issue:
\href{https://github.com/davislenover/BurnNEarn/issues/30}{\#30} (September 11),
\href{https://github.com/davislenover/BurnNEarn/issues/31}{\#31} (September 17) and
\href{https://github.com/davislenover/BurnNEarn/issues/257}{\#257} (November 5).
Only meetings with an issue are counted. The whole team also met at the four
TA discussions in Section 5, the first of which was also the supervisor
meeting in Section 3.
Day-to-day discussion is on Microsoft Teams and in GitHub issues and pull
requests (Development Plan, Section 5), which this table does not count.

\section{Supervisor/Stakeholder Meeting Attendance}

\noindent \textbf{Supervisor's Name: } Chris Schankula

\begin{table}[H]
\centering
\begin{tabular}{ll}
\toprule
\textbf{Student} & \textbf{Meetings}\\
\midrule
Total & 1\\
Davis Lenover & 1\\
Preston Grewal & 1\\
Angad Chhabra & 1\\
Farid Bastoros & 1\\
Zaid Alsaadi & 1\\
\bottomrule
\end{tabular}
\end{table}

The count comes from the supervisor meeting issue
\href{https://github.com/davislenover/BurnNEarn/issues/33}{\#33} (September 23).
Chris Schankula is also the team's TA and one of its stakeholders. That
meeting was also the first TA document discussion, so it is counted in
Section 5 as well. The team met him again at the three later discussions,
which are counted only in Section 5. At the last one, on November 11, the team
also showed him the proposed app screens for his feedback as a stakeholder
(\href{https://github.com/davislenover/BurnNEarn/issues/311}{\#311}).
\href{https://github.com/davislenover/BurnNEarn/issues/34}{\#34} is an unused
copy of the meeting template and is not counted.

\section{Lecture Attendance}

\begin{table}[H]
\centering
\begin{tabular}{ll}
\toprule
\textbf{Student} & \textbf{Lectures}\\
\midrule
Total & 6\\
Davis Lenover & 6\\
Preston Grewal & 0\\
Angad Chhabra & 2\\
Farid Bastoros & 5\\
Zaid Alsaadi & 1\\
\bottomrule
\end{tabular}
\end{table}

The total is the number of lecture issues the team created, and each member's
count comes from the attendance lists in those issues:

\begin{itemize}
\item \href{https://github.com/davislenover/BurnNEarn/issues/1}{\#1} (September 21, POC and Development Plan): Davis, Angad and Farid
\item \href{https://github.com/davislenover/BurnNEarn/issues/35}{\#35} (September 24, Requirements): Davis
\item \href{https://github.com/davislenover/BurnNEarn/issues/53}{\#53} (September 28, Requirements): Davis, Angad and Farid
\item \href{https://github.com/davislenover/BurnNEarn/issues/73}{\#73} (October 5, Hazard Analysis): Davis, Farid and Zaid
\item \href{https://github.com/davislenover/BurnNEarn/issues/145}{\#145} (October 21, VnV Plan and POC): Davis and Farid
\item \href{https://github.com/davislenover/BurnNEarn/issues/217}{\#217} (November 2, Design Document): Davis and Farid
\end{itemize}

The course calendar lists 11 lectures between the start of classes and the POC
demos. The other five have no lecture issue, so they are not counted. Four of
them, on September 9, 10, 16 and 17, were held before the team opened its
first issues on September 21. The fifth is a guest lecture on September 23.

\section{TA Document Discussion Attendance}

\noindent \textbf{TA's Name: } Chris Schankula

\begin{table}[H]
\centering
\begin{tabular}{ll}
\toprule
\textbf{Student} & \textbf{Meetings}\\
\midrule
Total & 4\\
Davis Lenover & 4\\
Preston Grewal & 4\\
Angad Chhabra & 4\\
Farid Bastoros & 4\\
Zaid Alsaadi & 4\\
\bottomrule
\end{tabular}
\end{table}

The TA scheduled four document discussions with the team, and the counts come
from the attendance lists in their issues: \href{https://github.com/davislenover/BurnNEarn/issues/33}{\#33} (September 23, the
first informal review), \href{https://github.com/davislenover/BurnNEarn/issues/98}{\#98} (October 7, the SRS and the Hazard
Analysis), \href{https://github.com/davislenover/BurnNEarn/issues/184}{\#184} (October 28, the VnV Plan and POC planning) and
\href{https://github.com/davislenover/BurnNEarn/issues/311}{\#311} (November 11, the Design Document Revision $-1$ and POC plans).
All five members attended each one. \href{https://github.com/davislenover/BurnNEarn/issues/33}{\#33} is titled as a supervisor
meeting, since Chris Schankula is both the team's TA and its supervisor, so
that one meeting is counted here and in Section 3.

\section{Commits, Lines Added and Lines Deleted, Pull Requests}

\begin{table}[H]
\centering
\begin{tabular}{lllllll}
\toprule
\textbf{Student} & \textbf{Commits} & \textbf{Percent} & \textbf{Lines Added} &
\textbf{Lines Deleted} & \textbf{Pull Requests}\\
\midrule
Total & 666 & 100\% & 56,518 & 11,391 & 110\\
Davis Lenover & 300 & 45.0\% & 23,779 & 5,568 & 39\\
Preston Grewal & 60 & 9.0\% & 3,858 & 526 & 12\\
Angad Chhabra & 37 & 5.6\% & 13,515 & 2,809 & 9\\
Farid Bastoros & 199 & 29.9\% & 8,450 & 2,179 & 36\\
Zaid Alsaadi & 70 & 10.5\% & 6,916 & 309 & 14\\
\bottomrule
\end{tabular}
\end{table}

These numbers were pulled on November 19, 2026, at 8:00 PM, at commit
\texttt{a9f1cc6} of \texttt{main}, with \texttt{git log --no-merges --numstat}.
They will likely change before the POC demo, since the team is still making
adjustments to the POC before presenting it. The team merges
pull requests with merge commits, so the commits of a merged branch are on
\texttt{main} and count for their author. The percentages are rounded to one
decimal place.

Three kinds of commit are left out: the merge commits themselves, the 96
commits the CI bot made to rebuild the PDFs, and three commits that only bring
in or remove the course template's text. Those three are the initial commit,
which created the repository from the course template, the commit that added
the course's IEEE SRS template, and the commit that removed the unused SRS
templates. Text a member wrote over the template is counted, since only the
lines a commit adds count for its author.

Binary files, such as the PDFs and the figures, add no lines. The lines added
include two generated lock files, 16,992 lines of the app's
\texttt{package-lock.json} and 1,081 lines of the server's \texttt{uv.lock},
which count for the member who committed them. By kind of file, 23,311 of the
lines added are in the \LaTeX{} documents, 18,073 in the lock files, 5,693 in
backend tests, 3,798 in backend code, 2,164 in figure sources, 1,567 in
Markdown notes, 995 in frontend code, 712 in backend migrations and 205 in CI
and other files.
Section 11 gives the same split for each member, with each member's net lines.

Line counts run high for documents. They are counted by source line, and most
of their text is wrapped at about 80 characters, so one paragraph is six to
eight lines and each table row is a line. A line that is revised counts again
as a line added, and the old one as a line deleted. So does a line that is
only indented differently or pasted again. Net of the lines deleted, the team
added 45,127 lines. These numbers show the size of what each member
committed, not how hard it was to write.

The pull requests are all those each member opened: 104 are merged, four were
closed without merging and two are open. The open ones are
\href{https://github.com/davislenover/BurnNEarn/pull/312}{\#312}, the POC demo
screen, and
\href{https://github.com/davislenover/BurnNEarn/pull/348}{\#348}, which adds
these sections to this report. Their commits are not in the counts above yet.
When these were counted, \#312 had 26 commits by Zaid Alsaadi, which add 688
lines of frontend code and delete 152. Section 11 shows the lines added with
that pull request included. No commit has more than one author.

\section{Issue Tracker}

% \wss{For each team member how many issues have they authored (including open and
% closed issues (O+C)) and how many have they been assigned (only counting closed
% issues (C only)) over the time period of interest.}

\begin{table}[H]
\centering
\begin{tabular}{lll}
\toprule
\textbf{Student} & \textbf{Authored (O+C)} & \textbf{Assigned (C only)}\\
\midrule
Davis Lenover & 70 & 55 \\
Preston Grewal & 1 & 40 \\
Angad Chhabra & 1 & 37 \\
Farid Bastoros & 119 & 47 \\
Zaid Alsaadi & 7 & 37 \\
\bottomrule
\end{tabular}
\end{table}

These counts cover the same period as Sections 2 to 6, up to November 18,
2026, and were pulled with GitHub's issue search. Authored counts every issue
a member opened in that period, open or closed. Assigned counts the issues
assigned to a member that were closed by November 18. Pull requests are not
counted, since Section 6 covers them. An issue assigned to more than one member
counts for each of them.



% \wss{If needed, an explanation for the counts can be provided here.}

\section{CICD via GitHub Actions}

\wss{Say how CICD technology is used and will be used in your project.}

\wss{Provide links to your CICD yaml files}

\section{Custom Documents (CDs)}

% \wss{What is the plan (as documented in TeamComposition.csv) for the team's
% CDs?  Should the CDs be modified now that the team knows more about the
% project?}

\paragraph{Plan.} The team's two CDs, one per term, are the ones in
TeamComposition.csv and in Section 4 of the Problem Statement and Goals:
Design Thinking in the Fall and Usability Testing in the Winter. The V\&V
Plan (Section 2.3) explains how both feed into validation.

\paragraph{Design Thinking (Fall): done.} The team followed the course's
Design Thinking process up to the second prototype, and the completed
35-slide worksheet is in the repository as
\href{https://github.com/davislenover/BurnNEarn/blob/main/docs/CDs/DesignThinking/DesignThinkingWorksheet.pdf}{PDF}
and \texttt{.pptx}, merged on November 18
(\href{https://github.com/davislenover/BurnNEarn/pull/333}{\#333},
\href{https://github.com/davislenover/BurnNEarn/pull/336}{\#336}). It covers
domain research on fitness tracking and beginner investing, two interviews
held on October 14, 2026 (P01, our supervisor and TA, and P02, a financial
services representative at a credit union), and two prototypes with user
feedback: a seven-screen Figma concept, and the POC demo screen. The
interviews are the source of most of the stakeholder requirements in SRS
Section 3 and of Personas 1 and 2 in SRS Section 4.

\paragraph{Usability Testing (Winter): no change needed.} Now that the team
knows more about the project, we still think usability testing is the right
CD for the Winter term. It tests the core loop that the POC shows, and it
provides the evidence for NFR2 (Usability), NFR8 (Accessibility) and goals G1
and G2. Its participant profile, tasks and survey are already set in the V\&V
Plan (Sections 3.7 and 6.2), and testing will start after the Rev 0 demo, once
there is a working app. Design Thinking has refined what the tests should
check:

\begin{itemize}
\item Neither interviewee was a beginner at investing, so participants will be
  recruited mainly from beginners, who are represented only by Persona 3.
\item Testers of Prototype 2 asked whether credits are the same as dollars and
  what happens after 300 credits. The tests will check whether users
  understand that credits are virtual, and where their credits came from.
\item When steps could not be read, the screen only showed ``Unavailable'', so
  the tests will include that case and check whether users understand why.
\end{itemize}

\section{Team Charter Trigger Items}

The Team Charter in the
\href{https://github.com/davislenover/BurnNEarn/blob/main/docs/DevelopmentPlan/DevelopmentPlan.tex}{Development Plan}
sets five triggers.

\begin{itemize}
  \item \textbf{Meeting attendance.} Attend fewer than 80\% of team meetings
    in a deliverable period, or miss two in a row without an excuse.
  \item \textbf{Missed team deadlines.} Miss two or more team deadlines in a term.
  \item \textbf{Review turnaround.} Take more than 72 hours from the review
    request to complete a review, two or more times in a deliverable period.
  \item \textbf{Responsiveness.} Miss the deadline to acknowledge a GitHub or
    Teams mention three or more times in a deliverable period. The deadline
    is 12 hours, or 24 hours on weekends.
  \item \textbf{Contribution.} Go two consecutive weeks without closing an
    assigned issue, contributing a merged pull request, or submitting a
    pull request review.
\end{itemize}

\begin{table}[H]
\caption{Team charter findings through November 18, 2026}
\centering
\begin{tabularx}{\linewidth}{@{}>{\raggedright\arraybackslash}p{0.22\linewidth}
  >{\raggedright\arraybackslash}X@{}}
\toprule
\textbf{Trigger} & \textbf{Finding}\\
\midrule
Meeting attendance &
  Three members attended 3 of 3 recorded meetings and two attended 2 of 3
  (67\%). No member missed two recorded meetings in a row, absences were excused.
  (\href{https://github.com/davislenover/BurnNEarn/issues/30}{\#30},
  \href{https://github.com/davislenover/BurnNEarn/issues/31}{\#31},
  \href{https://github.com/davislenover/BurnNEarn/issues/257}{\#257})\\
\addlinespace
Missed team deadlines & N/A\\
\addlinespace
Review turnaround & N/A\\
\addlinespace
Responsiveness & N/A\\
\addlinespace
Contribution & N/A\\
\bottomrule
\end{tabularx}
\end{table}

\paragraph{Follow-up:}
No violations currently to address.

\section{Additional Productivity Metrics}

% \wss{If your team has additional metrics of productivity, please feel free to
% add them to this report.  If not, please explicitly state that there are no additional
% metrics.}

% \wss{Additional metrics can include things like code reviews done, pull requests
% created, count of joining meetings late, count of number of times contributions
% had to be corrected, number of internal deadlines missed, test cases written, etc.}

% \wss{We are looking for data on these metrics, not just a list of additional
% metrics the team is planning on tracking.  However, if all you have is a plan,
% please share it here.}

\subsection{Pull Request Reviews and Merges}

Pull requests by other members that each member reviewed on GitHub, and pull
requests each member merged into \texttt{main}, split into another member's
and their own, up to November 18 (commit \texttt{2f3ab0c}).

\begin{table}[H]
\centering
\begin{tabular}{lrrr}
\toprule
\textbf{Student} & \textbf{Reviewed} & \textbf{Merged, another's} &
\textbf{Merged, own}\\
\midrule
Total & 93 & 82 & 15\\
Davis Lenover & 45 & 46 & 8\\
Preston Grewal & 5 & 0 & 0\\
Angad Chhabra & 5 & 0 & 0\\
Farid Bastoros & 36 & 34 & 7\\
Zaid Alsaadi & 2 & 2 & 0\\
\bottomrule
\end{tabular}
\end{table}

\subsection{Lines Added by Kind of File}

Section 6 gives each member's lines added. This table splits them by the kind
of file they are in, counted the same way and at the same commit. Like the
numbers in Section 6, these were pulled on November 19, 2026, at 8:00 PM and
will likely change before the POC demo.

\begin{table}[H]
\centering
% Smaller, with less space between columns, so the table fits the text width.
\small
\setlength{\tabcolsep}{3pt}
\begin{tabular}{lrrrrrr}
\toprule
\textbf{Kind of file} & \textbf{Davis} & \textbf{Preston} & \textbf{Angad} &
\textbf{Farid} & \textbf{Zaid} & \textbf{Team}\\
\midrule
Documents (\LaTeX{}) & 10,780 & 1,770 & 2,535 & 7,695 & 531 & 23,311\\
Generated lock files & 977 & 196 & 10,633 & 104 & 6,163 & 18,073\\
Backend tests (POC) & 4,635 & 806 & 0 & 252 & 0 & 5,693\\
Backend code (POC) & 3,252 & 302 & 0 & 244 & 0 & 3,798\\
Figure sources & 1,904 & 116 & 0 & 144 & 0 & 2,164\\
Markdown notes & 1,375 & 186 & 0 & 6 & 0 & 1,567\\
Backend migrations (POC) & 656 & 56 & 0 & 0 & 0 & 712\\
Frontend code (POC) & 0 & 426 & 347 & 0 & 222 & 995\\
CI and other files & 200 & 0 & 0 & 5 & 0 & 205\\
\midrule
All lines added & 23,779 & 3,858 & 13,515 & 8,450 & 6,916 & 56,518\\
Lines deleted & 5,568 & 526 & 2,809 & 2,179 & 309 & 11,391\\
Net lines (added less deleted) & 18,211 & 3,332 & 10,706 & 6,271 & 6,607 & 45,127\\
\midrule
Open PR \#312, frontend (POC) & 0 & 0 & 0 & 0 & 688 & 688\\
Lines added with the open PR & 23,779 & 3,858 & 13,515 & 8,450 & 7,604 & 57,206\\
\bottomrule
\end{tabular}
\end{table}

The backend rows are the POC server: the files under \texttt{src/backend}
(its code, README and configuration, apart from its lock file), its database
migrations, and its tests under \texttt{test/backend}. Frontend code
is the POC app under \texttt{src/frontend}, with its test scripts, apart from
its lock file. The POC
demo screen is still in the open pull request
\href{https://github.com/davislenover/BurnNEarn/pull/312}{\#312}, so it is not
on \texttt{main} and not in the frontend row. The row under the total adds its
26 commits, counted the same way, and the last row gives each member's lines
added with it included. Figure sources are the PlantUML and DOT files the diagrams are drawn from.
The drawn figures and the PDFs are binary files, which git counts as no lines.
One drawn SVG, since replaced by a PNG, adds two lines to the figure sources
row. Markdown notes are the architecture notes, \texttt{CONTRIBUTING.md} and the README
files outside \texttt{src/}.

\end{document}